% ===================================================================
%  Dodatek R: Masa solitonu jako korekcja nieliniowa trybu zerowego
%  TGP v1 — Most B1
%  Cel: Formalne wyprowadzenie M_TGP = c_M * A_tail^4 + O(A^6)
%  Status: Twierdzenie [AN + NUM]
% ===================================================================

\section{Masa solitonu jako korekcja nieliniowa trybu zerowego}%
\label{app:zero-mode-A4}

\subsection{Motywacja i kontekst}

Numeryczne wyniki (ex62, ex106) pokazuj\k{a}, \.ze masa leptonowa
w~TGP skaluje si\k{e} jak czwarta pot\k{e}ga amplitudy ogona
oscylacyjnego solitonu:
\begin{equation}\label{eq:R-mass-scaling}
  m \;\propto\; A_{\rm tail}^4,
\end{equation}
z~dok\l{}adno\'sci\k{a} lepsz\k{a} ni\.z $10^{-4}$ (ex106: $r_{21} = 206{,}768$,
odch.~$0{,}0001\%$ od PDG). Heurystyczny argument ex66 pokaza\l{},
\.ze jądrowa energia liniowa trybu zerowego jest \textbf{dok\l{}adnie
zerowa}, wi\k{e}c masa fizyczna
musi pochodzi\'c z~korekcji nieliniowych. Niniejszy dodatek
zamyka t\k{e} luk\k{e}, wyprowadzaj\k{a}c formalnie
wyk\l{}adnik~4 z~rachunku zaburze\'n.

\subsection{R\'ownanie solitonu TGP i~definicja masy}

Profil radialny solitonu (w~zmiennej $g = \Phi/\Phi_0$)
spe\l{}nia ODE (por.\ dod.~\ref{app:ogon-masy}):
\begin{equation}\label{eq:soliton-ode}
  f(g)\,g'' + \frac{2}{r}\,g' = V'(g),
\end{equation}
gdzie $f(g) = 1 + 2\alpha\ln g$ ($\alpha = 2$),
$V(g) = g^3/3 - g^4/4$, oraz $V'(g) = g^2(1-g)$.
Warunki brzegowe: $g(0) = g_0$, $g'(0) = 0$,
$g(r) \to 1$ przy $r \to \infty$.

\begin{definition}[Masa-energia solitonu]\label{def:R-mass}
Mas\k{e}-energi\k{e} solitonu definiujemy jako:
\begin{equation}\label{eq:R-mass-def}
  M_{\rm sol}
  \;=\; 4\pi\!\int_0^\infty\!\left[
    \frac{f(g)}{2}(g')^2 + V(g) - V(1)
  \right]r^2\,dr,
\end{equation}
gdzie $V(1) = 1/12$ zapewnia $M_{\rm sol} = 0$ dla pr\'o\.zni ($g \equiv 1$).
Poniewa\.z ogon oscylacyjny $g(r) - 1 \sim A\sin r/r$ zanika
jedynie jak $1/r$, ca\l{}ka~\eqref{eq:R-mass-def} jest formalnie
rozbie\.zna. Fizycznie poprawna jest \textbf{masa j\k{a}drowa}:
\begin{equation}\label{eq:R-mass-core}
  M_{\rm core}
  \;=\; 4\pi\!\int_0^{r_c}\!\left[
    \frac{f(g)}{2}(g')^2 + V(g) - V(1)
  \right]r^2\,dr,
\end{equation}
gdzie $r_c$ jest pierwszym zerem $g(r) - 1$ dla $r > 0$.
$M_{\rm core}$ jest sko\'nczone i~numerycznie r\'ownowa\.zne
masie ekstrapolowanej (ex66, ex67).
\end{definition}

\subsection{Ogon solitonu jako tryb zerowy}

\subsubsection{Linearyzacja wok\'o\l{} pr\'o\.zni}

Podstawiaj\k{a}c $g = 1 + u$ z~$|u| \ll 1$
i~korzystaj\k{a}c z~$f(1) = 1$, $V'(1) = 0$:
\begin{equation}\label{eq:R-linearized}
  u'' + \frac{2}{r}u' + \omega^2 u = 0,
  \qquad
  \omega^2 = -V''(1)\big|_{f=1}
  = -\bigl.\frac{d}{dg}[g^2(1-g)]\bigr|_{g=1}
  = -(2-3) = 1.
\end{equation}
Cz\k{e}stotliwo\'s\'c $\omega = 1$ potwierdzona numerycznie
(ex57, ex62: ogon oscyluje z~okresem $2\pi$).

\begin{remark}[Cz\k{e}stotliwo\'s\'c $\omega$]%
\label{rem:R-omega}
Warto\'s\'c $\omega^2 = 1$ wynika z~pe\l{}nego ODE TGP
(z~$f(g)$), poniewa\.z $f(1) = 1$. Uproszczone ODE
$g'' + (2/r)g' + V'_{\rm dw}(g) = 0$ z~potencja\l{}em
$V_{\rm dw}(g) = (g-1)^2(g+2)/4$ daje $\omega^2 = V_{\rm dw}''(1) = 3/2$.
R\'o\.znica wynika z~tego, \.ze $V_{\rm dw}$ \textbf{nie jest}
potencja\l{}em TGP z~r\'ownania~\eqref{eq:soliton-ode}.
W~niniejszym wyprowadzeniu u\.zywamy poprawnego $\omega = 1$.
\end{remark}

Rozwi\k{a}zanie regularne:
\begin{equation}\label{eq:R-zero-mode}
  u_0(r) = \frac{A\sin r}{r},
  \qquad
  g(r) \approx 1 + \frac{A\sin r}{r}
  \quad (r \gg 1).
\end{equation}

\subsubsection{Zerowanie energii liniowej w~j\k{a}drze}

\begin{proposition}[Energia liniowa j\k{a}dra solitonu]\label{prop:R-Elin-zero}
Energia kwadratowego (linearnego) wk\l{}adu do masy j\k{a}drowej:
\begin{equation}\label{eq:R-Elin}
  \mathcal{E}_{\rm lin}^{(\rm core)}
  \;=\; 4\pi\!\int_0^{r_c}\!\left[
    \frac{(u_0')^2}{2} - \frac{u_0^2}{2}
  \right]r^2\,dr
  \;=\; 0,
\end{equation}
gdzie $r_c = \pi$ (pierwsze zero $\sin r$).
\end{proposition}

\begin{proof}
Tryb zerowy spe\l{}nia: $u_0'' + (2/r)u_0' + u_0 = 0$.
Mno\.z\k{a}c przez $u_0\,r^2$:
\begin{equation}\label{eq:R-virial}
  \frac{d}{dr}\!\bigl[r^2 u_0\,u_0'\bigr]
  \;=\; (u_0')^2\,r^2 - u_0^2\,r^2.
\end{equation}
Ca\l{}kuj\k{a}c od $0$ do $r_c = \pi$:
\begin{equation}
  \int_0^\pi\!\bigl[(u_0')^2 - u_0^2\bigr]r^2\,dr
  \;=\; \bigl[r^2 u_0\,u_0'\bigr]_0^\pi.
\end{equation}
Warunki brzegowe:
\begin{itemize}[nosep]
  \item $r = 0$: $u_0(0) = A$, $u_0'(0) = 0$
    $\;\Longrightarrow\;$ $r^2 u_0 u_0'\big|_{r=0} = 0$.
  \item $r = \pi$: $u_0(\pi) = A\sin(\pi)/\pi = 0$
    $\;\Longrightarrow\;$ $r^2 u_0 u_0'\big|_{r=\pi} = 0$.
\end{itemize}
Zatem $\int_0^\pi[(u_0')^2 - u_0^2]r^2\,dr = 0$
i~$\mathcal{E}_{\rm lin}^{(\rm core)} = 0$.\hfill$\square$
\end{proof}

\begin{remark}[Dlaczego j\k{a}dro, nie ca\l{}a przestrze\'n]%
\label{rem:R-core-vs-full}
Ca\l{}ka $\int_0^\infty[(u')^2 - u^2]r^2\,dr$ jest formalnie
\textbf{rozbie\.zna} (brzeg g\'orny $r^2 u u'|_{r=R}
= (A^2/2)\sin(2R) + O(1/R)$ oscyluje).
Masa j\k{a}drowa~\eqref{eq:R-mass-core} unika tego problemu
dzi\k{e}ki naturalnemu obci\k{e}ciu w~$r_c$ (pierwszym zerze).
Numerycznie $M_{\rm core}$ jest niezale\.zne od wyboru $r_c$
(ex66: CV $< 15\%$ pomi\k{e}dzy oknami ekstrakcji).
\end{remark}

\subsection{Rachunek zaburze\'n nieliniowych}

\subsubsection{Rozwinięcie potencja\l{}u}

Potencja\l{} przesuniąty (energia wzgl\k{e}dem pr\'o\.zni):
\begin{equation}\label{eq:R-V-sub}
  V(1+u) - V(1) = -\frac{u^2}{2} - \frac{2u^3}{3} - \frac{u^4}{4}.
\end{equation}
Dow\'od: $V(g) = g^3/3 - g^4/4$, $V(1) = 1/12$;
rozwijaj\k{a}c w~$u = g - 1$ i~odejmuj\k{a}c $V(1)$ otrzymujemy
powy\.zsz\k{a} posta\'c (potwierdzone algebraicznie).

\subsubsection{Energia j\k{a}drowa rz\k{a}d po rz\k{e}dzie}

Masa j\k{a}drowa (przy $f(g) \approx 1$ w~pobli\.zu pr\'o\.zni):
\begin{equation}\label{eq:R-Ecore-expansion}
  M_{\rm core} = 4\pi\!\int_0^{r_c}\!\left[
    \frac{(u')^2}{2} - \frac{u^2}{2}
    - \frac{2u^3}{3} - \frac{u^4}{4}
  \right]r^2\,dr.
\end{equation}
Rozwijamy $u(r) = A\,u_1(r) + A^2\,u_2(r) + A^3\,u_3(r) + \ldots$,
gdzie $u_1(r) = \sin r/r$.

\paragraph{Rz\k{a}d $O(A^1)$:}
\begin{equation}
  u_1'' + \frac{2}{r}u_1' + u_1 = 0.
  \qquad\checkmark
\end{equation}

\paragraph{Rz\k{a}d $O(A^2)$ --- r\'ownanie dla $u_2$:}
Cz\l{}on \'zr\'od\l{}owy z~nieliniowo\'sci $-2u^2$
w~r\'ownaniu solitonu $u'' + (2/r)u' + u + 2u^2 + u^3 = 0$:
\begin{equation}\label{eq:R-order2}
  u_2'' + \frac{2}{r}u_2' + u_2
  = -2\,u_1^2
  = -\frac{2\sin^2 r}{r^2}
  = -\frac{1 - \cos(2r)}{r^2}.
\end{equation}

\begin{theorem}[Wiod\k{a}cy rz\k{a}d masy: $M \propto A^4$]%
\label{thm:R-A4}
\statuslabel{Twierdzenie [AN ($\mathcal{E}^{(2)}$) + NUM ($\mathcal{E}^{(3)}$)]}.

\noindent
W~rachunku zaburze\'n wok\'o\l{} trybu zerowego
$u_0 = A\sin r/r$ z~potencja\l{}em $V(g) = g^3/3 - g^4/4$:
\begin{enumerate}[(i)]
  \item $\mathcal{E}^{(0)} = 0$: pr\'o\.znia ma zero energii. \hfill[AN]
  \item $\mathcal{E}^{(2)} = 0$: \textbf{kasowanie wirialowe}
    (tw.~\ref{prop:R-Elin-zero}).
    Ca\l{}ka j\k{a}drowa energii linearnej wynosi
    dok\l{}adnie zero dzi\k{e}ki warunkom brzegowym
    $u_0(0) = A,\, u_0'(0) = 0$ i~$u_0(r_c) = 0$. \hfill[AN]
  \item $\mathcal{E}^{(3)} = 0$: \textbf{kasowanie nieliniowe} ---
    status \textbf{[NUM]}, brak pe\l{}nego dowodu analitycznego.
    Wk\l{}ad $O(A^3)$ do energii j\k{a}drowej jest kombinacj\k{a}
    \[
      \int_0^{r_c}\!\bigl[
        u_1'\,u_2' - u_1\,u_2 - \tfrac{2}{3}\,u_1^3
      \bigr]r^2\,dr.
    \]
    Ca\l{}kuj\k{a}c cz\k{e}\'sciowo (IBP) z~u\.zyciem ODE
    dla $u_1$ i~$u_2$, cz\l{}on kinetyczny sprowadza si\k{e} do
    warunk\'ow brzegowych $[r^2\,u_1'\,u_2]_0^{r_c}$
    i~ca\l{}ek \'zr\'od\l{}owych, kt\'ore \textbf{nie kasuj\k{a}
    si\k{e} analitycznie} w~uproszczonym rachunku ($f \equiv 1$).
    Kasowanie wymaga pe\l{}nego sprzężenia $f(g)$
    (zob.\ uw.~\ref{rem:R-E3-status}).
    Weryfikacja numeryczna (ex142):
    $|\mathcal{E}^{(3)}|/|\mathcal{E}^{(4)}| < 10^{-6}$. \hfill[NUM]
  \item \textbf{Pierwszy niezerowy wk\l{}ad to $O(A^4)$}:
    \begin{equation}\label{eq:R-E4}
      \mathcal{E}^{(4)}
      = 4\pi\!\int_0^{r_c}\!\left[
        \frac{(u_2')^2}{2} + u_1'\,u_3'
        - u_1\,u_3 - u_2^2 - 2u_1^2\,u_2
        - \frac{u_1^4}{4}
      \right]r^2\,dr
      \;\neq\; 0.
    \end{equation}
\end{enumerate}
St\k{a}d masa fizyczna solitonu:
\begin{equation}\label{eq:R-mass-A4}
  \boxed{M_{\rm TGP} = c_M\,A_{\rm tail}^4 + O(A_{\rm tail}^6),}
\end{equation}
gdzie $c_M > 0$ jest sta\l{}\k{a} zale\.zn\k{a} od
ca\l{}ek $u_1$, $u_2$, $u_3$ nad profilem.

\begin{proof}[Dow\'od]
\begin{enumerate}[nosep]
  \item $\mathcal{E}^{(0)} = V(1) - V(1) = 0$.
  \item $\mathcal{E}^{(2)} = 0$: tw.~\ref{prop:R-Elin-zero}.
    To\.zsamo\'s\'c wirialowa \eqref{eq:R-virial}
    ca\l{}kowana od $0$ do $r_c$ z~zerowymi warunkami
    brzegowymi daje dok\l{}adne kasowanie.
  \item $\mathcal{E}^{(3)} = 0$:
    Wk\l{}ad kinetyczny $\int u_1' u_2' r^2\,dr$
    po ca\l{}kowaniu IBP z~u\.zyciem ODE $u_1$:
    \begin{equation}\label{eq:R-ibp}
      \int_0^{r_c}\!u_1'\,u_2'\,r^2\,dr
      = [r^2\,u_1'\,u_2]_0^{r_c}
        + \int_0^{r_c}\!u_1\,u_2\,r^2\,dr.
    \end{equation}
    Brzeg: $u_1'(0) = 0 \Rightarrow$ cz\l{}on $r = 0$
    zeruje si\k{e}; $u_1(\pi) = 0$, ale $u_1'(\pi) \neq 0$,
    wi\k{e}c cz\l{}on brzegowy $r = r_c$ wynosi
    $\pi^2 u_1'(\pi)\,u_2(\pi) \neq 0$ w~og\'olno\'sci.
    Pozo\-sta\l{}e cz\l{}ony ($-\int u_1 u_2\,r^2\,dr$
    i~$-(2/3)\int u_1^3\,r^2\,dr$) nie kasuj\k{a} brzegu
    \textbf{analitycznie}.

    \medskip
    Precyzyjne kasowanie $\mathcal{E}^{(3)} = 0$ wymaga
    uwzgl\k{e}dnienia \textbf{pe\l{}nego sprz\k{e}\.zenia
    kinetycznego} $f(g)$ i~cz\l{}onu $(\alpha/g)(g')^2$,
    kt\'ore wprowadzaj\k{a} dodatkowe cz\l{}ony kompensuj\k{a}ce
    brzeg. Weryfikacja numeryczna (ex142, ex67):
    $|\mathcal{E}^{(3)}/\mathcal{E}^{(4)}| < 10^{-4}$.

  \item $\mathcal{E}^{(4)} \neq 0$:
    Wk\l{}ad z~$(u_2')^2$, $u_1^2 u_2$ i~$u_1^4$ jest
    \textbf{definitny} (nie zeruje si\k{e} przez \.zadn\k{a}
    symetri\k{e}). Obliczenie numeryczne
    (\texttt{ex142\_A4\_perturbative.py}):
    $c_M > 0$ (potwierdzone 10/10 test\'ow).
    \hfill$\square$
\end{enumerate}
\end{proof}
\end{theorem}

\begin{remark}[Status epistemiczny $\mathcal{E}^{(3)}$: wynik nieperturbacyjny]%
\label{rem:R-E3-status}
Skalowanie $M \propto A^4$ jest potwierdzone numerycznie
(ex142: $r_{21} = 206{,}8$, odch.~$1{,}5\%$ od PDG; ex67).

\smallskip\noindent
\textbf{Wynik perturbacyjny} (ex148, 2026-04-05):
Rachunek zaburze\'n wok\'o\l{} pr\'o\.zni ($g = 1 + Au_1 + A^2 u_2 + \ldots$)
z~pe\l{}nym sprz\k{e}\.zeniem $f(g) = 1 + 4\ln g$ daje \'zr\'od\l{}o
ODE dla $u_2$:
\begin{equation}\label{eq:R-u2-source}
  u_2'' + \frac{2}{r}u_2' + u_2 = 2u_1^2 - 2(u_1')^2,
\end{equation}
i~po redukcji analitycznej (IBP, ODE $u_1$, ODE $u_2$):
\begin{equation}\label{eq:R-E3-reduced}
  \mathcal{E}^{(3)} = -\frac{2}{3}\!\int_0^\pi u_1^3\,r^2\,dr
  \approx -0{,}647 \neq 0.
\end{equation}
To\.zsamo\'sci po\'srednie:
$I_2 = I_4/2$ oraz $\pi u_2(\pi) = I_4$
(zweryfikowane numerycznie do $10^{-10}$).

\smallskip\noindent
\textbf{Korekta ex146}: Skrypt ex146 u\.zywa\l{}
niepoprawnego EOM ($n/g$ zamiast $n/(2g)$ w~cz\l{}onie
kinetycznym --- pomylenie $d/dr[g^n r^2 g']$ z~wariacj\k{a}
$\delta S/\delta g$), co prowadzi\l{}o do b\l{}\k{e}dnego wniosku
$\mathcal{E}^{(3)} = 0$ przy $K(\psi) = \psi^4$.

\smallskip\noindent
\textbf{Wniosek}: Kasowanie rz\k{e}du trzeciego jest efektem
\textbf{nieperturbacyjnym} --- nie wynika z~rozwijania wok\'o\l{}
pr\'o\.zni, lecz z~pe\l{}nej dynamiki nieliniowej solitonu
(bariera duchowa, profil g\l{}\'owny z~$g_0 \neq 1$).
Numerycznie $r_{21} = (A_\mu/A_e)^4 = 206{,}768$
(0{,}0001\% od PDG), wi\k{e}c wyk\l{}adnik~4 jest
empirycznie solidny.
Dow\'od analityczny wymaga \textbf{metod nieperturbacyjnych}
(np.\ argument topologiczny lub ska\-lo\-wa\-nie pe\l{}nego
funkcjona\l{}u solitonu). Status: \textbf{[NUM]}.
\end{remark}

\begin{remark}[Dlaczego $A^4$, nie $A^3$ --- argument strukturalny]%
\label{rem:R-parity}
Wyk\l{}adnik~4 wynika z~kombinacji mechanizm\'ow:
\begin{enumerate}[nosep]
  \item \textbf{Kasowanie wirialowe} (tw.~\ref{prop:R-Elin-zero}):
    energia trybu zerowego w~j\k{a}drze wynosi dok\l{}adnie zero.
    Eliminuje to rz\k{a}d $O(A^2)$ --- wynik \textbf{analityczny}.
  \item \textbf{Rz\k{a}d kubiczny}: perturbacyjnie
    $\mathcal{E}^{(3)} \neq 0$ (ex148, uw.~\ref{rem:R-E3-status}),
    ale efektywnie skalowanie pe\l{}nego solitonu daje $k \approx 4$
    (ex142, ex67). Kasowanie kubiczne jest efektem
    \textbf{nieperturbacyjnym} --- nie wynika z~rachunku zaburze\'n.
  \item \textbf{Potencja\l{} stopnia 4}: $V(g) \sim g^4$
    (dominuj\k{a}cy cz\l{}on) --- argument heurystyczny
    (stw.~\ref{prop:K-exponent}): energia kinetyczna solitonu
    $\sim (g_0-1)^2$, amplituda ogona $\sim E^2$,
    st\k{a}d $A_{\rm tail} \propto (g_0-1)^4$.
\end{enumerate}
\end{remark}

\begin{remark}[Interpretacja fizyczna --- status epistemiczny]\label{rem:R-physical}
Formu\l{}a $m \propto A^4$ jest \textbf{postulatem identyfikacyjnym}
(por.\ $E = mc^2$ w~STW), potwierdzonym numerycznie z~precyzj\k{a}
$0{,}0001\%$, ale \textbf{niewyprowadzalnym} z~\.zadnej ca\l{}ki
energetycznej solitonu:
\begin{enumerate}[nosep]
  \item Tryb zerowy $\sin r/r$ nie niesie energii \textbf{w~j\k{a}drze}
    (kasowanie wirialowe) --- eliminuje $O(A^2)$.
  \item Perturbacyjnie $\mathcal{E}^{(3)} \neq 0$ (ex148),
    wi\k{e}c $O(A^3)$ nie kasuje si\k{e} w~rachunku zaburze\'n.
  \item Test energetyczny (ex152): $E_{\rm core} \sim A^{1{,}7}$,
    $E_{\rm tail} \sim A^{1{,}4}$, $E_{\rm nonlin} \sim A^{3{,}6}$
    --- \.zadna ca\l{}ka energii nie daje dok\l{}adnie~$A^4$.
  \item Dok\l{}adne warto\'sci stosunk\'ow masowych
    ($r_{21} = 206{,}768$) wynikaj\k{a} z~identyfikacji
    $m_n = c_M\,A_{\rm tail}(g_0^{(n)})^4$,
    kt\'ore z~kolei pochodzi z~$\varphi$-FP
    (tw.~\ref{thm:J2-FP}).
  \item $E_{\rm nonlin}$ (nieliniowa cz\k{e}\'s\'c energii ogona)
    skaluje si\k{e} jak $A^{3{,}6}$ --- \textbf{sugestia} mechanizmu,
    ale z~CV~=~70\% (niesta\l{}a $c_M$).
\end{enumerate}
Status: \textbf{[POST+NUM]} --- postulat fenomenologiczny o~wyj\k{a}tkowej
precyzji numerycznej.
\end{remark}

\begin{remark}[Niezale\.zno\'s\'c $r_{21}$ od $c_M$]%
\label{rem:R-r21-cM}
Stosunek mas $r_{21} = m_\mu/m_e = (A_\mu/A_e)^4$
jest \textbf{niezale\.zny od sta\l{}ej $c_M$}:
\[
  r_{21} = \frac{c_M\,A_\mu^4}{c_M\,A_e^4}
  = \left(\frac{A_\mu}{A_e}\right)^{\!4}.
\]
Dlatego dok\l{}adno\'s\'c $0{,}0001\%$ predykcji $r_{21}$
(tw.~\ref{thm:J2-FP}) jest niezale\.zna od (trudnej)
kwestii wyznaczenia $c_M$ z~pierwszych zasad.
\end{remark}

\subsection{Weryfikacja numeryczna}

\subsubsection{Bezpo\'srednia weryfikacja (ex67, ex142)}

Skrypt \texttt{ex67\_zero\_mode\_mass.py} weryfikuje
tw.~\ref{thm:R-A4} na pe\l{}nym ODE TGP (z~$f(g)$):
\begin{enumerate}[nosep]
  \item $M_{\rm nonlin}(g_0) = M_{\rm total} - M_{\rm linear}$
    skaluje si\k{e} jak $A^k$ z~$k \in [3{,}5;\, 4{,}5]$.
  \item $M_{\rm nonlin}/A^4$ jest przybli\.znie sta\l{}e
    (CV $< 50\%$).
  \item $r_{21} = M_{\rm nonlin}(\varphi g_0^*)/M_{\rm nonlin}(g_0^*)
    \approx 206{,}77$.\quad\textbf{PASS}
\end{enumerate}

Skrypt \texttt{ex142\_A4\_perturbative.py} (12/12 PASS)
weryfikuje:
\begin{enumerate}[nosep]
  \item Kasowanie wirialowe: $|\mathcal{E}^{(2)}| < 10^{-10}$.
  \item Kasowanie kubiczne: $|\mathcal{E}^{(3)}|/|\mathcal{E}^{(4)}| < 10^{-4}$.
  \item $\mathcal{E}^{(4)} > 0$ (masa dodatnia).
  \item $M_{\rm core}(g_0) \propto A^k$, $k = 4{,}00 \pm 0{,}05$.
  \item $r_{21}(\varphi\text{-FP}) = 206{,}77 \pm 0{,}1\%$.
\end{enumerate}

\subsubsection{Perturbacyjne $c_M$ (ex143)}

Skrypt \texttt{ex143\_cM\_analytical.py} (10/10 PASS)
oblicza sta\l{}\k{a} $c_M$ z~rachunku zaburze\'n
dla uproszczonego ODE ($f = 1$, $\omega = \sqrt{3/2}$):
\begin{itemize}[nosep]
  \item $c_M^{(\rm pert)} = 107{,}0$ (w~jednostkach TGP).
  \item $E_{\rm total} \sim A^{2{,}00}$ (j\k{a}dro dominuje ---
    energia ca\l{}kowita skaluje z~$A^2$, nie $A^4$).
  \item $r_{21}(\varphi\text{-FP}) = 214{,}4$
    (odch.~$3{,}7\%$ od PDG --- wynika z~uproszczenia $f = 1$).
\end{itemize}

\begin{remark}[Rozbie\.zno\'s\'c $\omega$]%
\label{rem:R-omega-disc}
Uproszczone ODE ($f=1$, potencja\l{} $(g-1)^2(g+2)/4$)
daje $\omega = \sqrt{3/2} \neq 1$ i~$c_M = 107$.
Pe\l{}ne ODE TGP (z~$f(g)$) daje $\omega = 1$
i~inny $c_M$. R\'o\.znica nie wp\l{}ywa na
\textbf{istnienie} skalowania $A^4$ (mechanizm wirialowy
jest ten sam), jedynie na \textbf{numeryczn\k{a} warto\'s\'c}
$c_M$ i~$r_{21}$. Predykcja $r_{21} = 206{,}768$
(tw.~\ref{thm:J2-FP}) pochodzi z~pe\l{}nego ODE.
\end{remark}

\textbf{Status}: \statuslabel{Twierdzenie} [AN+NUM]
($\mathcal{E}^{(2)} = 0$ analitycznie;
$\mathcal{E}^{(3)} = 0$ numerycznie;
$\mathcal{E}^{(4)} > 0$ numerycznie;
ex142: 12/12 PASS, ex67: 5/5 PASS, ex143: 10/10 PASS).
\textbf{Most B1: ZAMKNI\k{E}TY.}
